%% Supplementary Information for Nature Computational Science Brief Communication
%% Heavy theory, metrics and computational workflow live here by design.

\documentclass[11pt,a4paper]{article}

\usepackage{amsmath,amssymb,amsthm,mathtools}
\usepackage{bm}
\usepackage{booktabs}
\usepackage{enumitem}
\usepackage{geometry}
\usepackage[colorlinks=true,linkcolor=blue,citecolor=blue,urlcolor=blue]{hyperref}
\geometry{margin=2.5cm}

\newtheorem{proposition}{Proposition}
\newtheorem{theorem}{Theorem}
\newtheorem{remark}{Remark}

\newcommand{\R}{\mathbb{R}}
\newcommand{\N}{\mathcal{N}}
\newcommand{\E}{\mathbb{E}}
\newcommand{\Haar}{\mu_{\mathrm{Haar}}}
\newcommand{\HAB}{H_{AB}}
\newcommand{\Xnew}{X^{\mathrm{n}}}
\newcommand{\Xold}{X^{\mathrm{o}}}
\newcommand{\pacc}{p_{\mathrm{acc}}}
\newcommand{\Dtheta}{D_\theta}
\newcommand{\ftheta}{f_\theta}
\newcommand{\that}{\hat t}
\newcommand{\xnoisy}{x^{\mathrm{noisy}}}
\newcommand{\xtilde}{\tilde x}
\newcommand{\cskip}{c_{\mathrm{skip}}}
\newcommand{\cout}{c_{\mathrm{out}}}
\newcommand{\cin}{c_{\mathrm{in}}}
\newcommand{\sigmaD}{\sigma_{\mathrm{data}}}

\title{\textbf{Supplementary Information}\\[4pt]
Path Sampling Diagnoses Protein Diffusion Models}
\author{Muhammad R.~Hasyim\\
\small Simons Center for Computational Physical Chemistry,
New York University, New York, NY~10003}
\date{\today}

\begin{document}
\maketitle
\tableofcontents
\newpage

\section{Supplementary Note 1: Path Probabilities and TPS Ratios}
\label{sec:path-prob}

\subsection{Extended state for Boltz-2 reverse diffusion}

Boltz-2 reverse diffusion is treated as an inhomogeneous discrete-time Markov
chain. Let $x_i \in \R^{M \times 3}$ be the atom-coordinate tensor at
denoising step $i$, and let $L$ be the number of denoising steps. The extended
state records every random input used by the sampler:
\begin{equation}
  \mathcal{S}_i = (x_i, R_i, \tau_i, \varepsilon_i, i),
\end{equation}
where $R_i \in \mathrm{SO}(3)$ is drawn from Haar measure, $\tau_i$ is a
translation drawn from $\N(0,s_{\mathrm{trans}}^2 I)$, and
$\varepsilon_i$ is Gaussian churn noise. The reverse step is deterministic
once these random variables are fixed.

For the MEK1--FZC case study, state A denotes the high-noise initial state and
state B denotes the terminal structured state produced by the sampler. More
selective B definitions, such as pLDDT or RMSD thresholds, can be substituted
without changing the path-probability derivation.

\subsection{One-step transition kernel}

Each denoising step consists of SE(3) augmentation, noise injection and a
deterministic EDM-preconditioned update~\cite{karras2022}. Written in compact
form,
\begin{align}
  \xtilde_i &= R_i (x_i-\bar x_i) + \tau_i,\\
  \xnoisy_i &= \xtilde_i + \varepsilon_i,\\
  x_{i+1} &= T_i(x_i;R_i,\tau_i,\varepsilon_i),
\end{align}
where $\bar x_i$ is the atom centroid and $T_i$ includes the Boltz-2 denoiser.
The corresponding kernel is
\begin{align}
  p(&x_{i+1}, R_i,\tau_i,\varepsilon_i \mid x_i) \nonumber\\
  &=
  \Haar(R_i)\,
  \N(\tau_i;0,s_{\mathrm{trans}}^2I)\,
  \N(\varepsilon_i;0,v_iI)\,
  \delta\!\left(x_{i+1}-T_i(x_i;R_i,\tau_i,\varepsilon_i)\right),
  \label{eq:transition-kernel}
\end{align}
with churn variance $v_i$ determined by the EDM noise schedule.

\subsection{Extended path density}

For a complete path
$\mathcal{X}_{\mathrm{ext}}=\{x_0,(R_i,\tau_i,\varepsilon_i)_{i=0}^{L-1},x_L\}$,
the density in random-variable coordinates is
\begin{equation}
  \log P_{\mathrm{ext}}(\mathcal{X}_{\mathrm{ext}})
  =
  \log \rho(x_0)
  +
  \sum_{i=0}^{L-1}
  \left[
    \log \Haar(R_i)
    + \log \N(\tau_i;0,s_{\mathrm{trans}}^2I)
    + \log \N(\varepsilon_i;0,v_iI)
  \right].
  \label{eq:log-path}
\end{equation}
This is the key simplification: no neural-network likelihood is needed, because
the coordinates are deterministic functions of explicit random draws.

\subsection{Jacobian diagnostics}

Equation~\eqref{eq:log-path} is the target density for paths expressed in the
extended random-variable parameterization. If one instead reports a
coordinate-density diagnostic for the map
$\varepsilon_i \mapsto x_{i+1}$, a change-of-variables term appears:
\begin{equation}
  -\log |\det J_i|,\qquad
  J_i = \frac{\partial x_{i+1}}{\partial \varepsilon_i}.
\end{equation}
For an EDM-preconditioned denoiser,
\begin{equation}
  J_i =
  \beta_i I
  +
  \mu_i \nabla_z \ftheta(z)\big|_{z=\cin(\that_i)\xnoisy_i},
\end{equation}
where $\beta_i=1+\alpha_i(1-\cskip(\that_i))$ and
$\mu_i=-\alpha_i\cout(\that_i)\cin(\that_i)$. The implementation may log a
scalar schedule-only approximation for diagnostics, but this approximation is
not a calibrated neural-map determinant and should not be mixed into a
Hastings ratio unless the proposal density is derived in the same coordinates.

\subsection{Reactive path ensemble}

The fixed-length reactive ensemble is
\begin{equation}
  \pi(\mathcal{X}) \propto
  H_A(x_0)\,P_{\mathrm{ext}}(\mathcal{X})\,H_B(x_L),
\end{equation}
where $H_A$ and $H_B$ are indicator functions for the endpoint state volumes.
The Monte Carlo chain samples this ensemble with three proposal classes.

\subsubsection*{Forward shooting}

Select a shooting index $k$, keep the prefix $x_0,\ldots,x_k$, draw fresh
$(R_i,\tau_i,\varepsilon_i)$ from the prior for $i\ge k$, and integrate
forward. The proposal density for the new suffix is exactly the same product
of priors appearing in the target density. Hence
\begin{equation}
  \Delta_{\mathrm{fwd}}=0,\qquad
  \pacc^{\mathrm{fwd}}=\HAB(\Xnew).
\end{equation}
The same cancellation applies to Haar and translation terms. Scalar
schedule-only Jacobians also cancel for fixed-length paths, because their
values depend only on $i$.

\subsubsection*{Backward shooting}

Backward shooting proposes a new prefix by approximately inverting the denoiser
from $x_{k+1}$ to $x_k$. The fixed-point inversion solves
\begin{equation}
  x_{i+1}=(1+\alpha_i)\xnoisy_i-\alpha_i\Dtheta(\xnoisy_i,\that_i).
\end{equation}
Because this proposal is not the time reversal of the forward kernel, the
Hastings correction is explicit:
\begin{equation}
  \Delta_{\mathrm{bwd}}
  =
  [\log p_{\mathrm{fwd}}(X'_{0:k})-\log p_{\mathrm{fwd}}(X_{0:k})]
  +
  [\log q_{\mathrm{bwd}}(X_{0:k})-\log q_{\mathrm{bwd}}(X'_{0:k})]
  +
  \log\rho(x'_0)-\log\rho(x_0).
  \label{eq:delta-bwd}
\end{equation}
The move is accepted with
$\min(1,\exp(\Delta_{\mathrm{bwd}}))$ if the endpoint tests are satisfied.

\subsubsection*{Global reshuffle}

Global reshuffling draws a fresh $x_0$ and all random variables from their
priors, then integrates the entire denoising trajectory. The proposal is the
target prior over complete paths, so the Metropolis ratio is again unity after
endpoint checks. This move decorrelates the chain and ensures that disconnected
reactive regions can be reached.

\begin{theorem}[Ergodic move set]
The combination of forward shooting, backward shooting and global reshuffling
is irreducible on the fixed-length reactive path ensemble, provided every
reactive path has nonzero prior probability under the sampler.
\end{theorem}
\begin{proof}
Forward and backward shooting update suffixes and prefixes while satisfying
Metropolis--Hastings balance. Global reshuffling can propose any complete
reactive path directly from the prior with nonzero probability. Thus every
reactive path communicates with every other through the reshuffle move, and
the endpoint-restricted chain is irreducible.
\end{proof}

\section{Supplementary Note 2: Quotient-Space Diffusion Context}
\label{sec:quotient}

Molecular structures are naturally defined modulo global rigid motions:
coordinates $x$ and $g\cdot x$ represent the same intrinsic structure for
$g\in \mathrm{SE}(3)$. Recent quotient-space diffusion theory formalizes this
idea by treating samples as elements of $\mathcal{M}/\mathcal{G}$ rather than
the ambient coordinate space~\cite{xu2026}. For a projection
$\pi:\mathcal{M}\rightarrow\mathcal{M}/\mathcal{G}$, directions tangent to the
group orbit are vertical degrees of freedom, while orthogonal directions form
a horizontal subspace that changes the intrinsic molecular shape.

The manuscript uses this perspective to specify the path state for the
implemented Boltz-2 sampler; it does not introduce a new quotient-space
sampler. Boltz-2's random rotations and translations are not physical
conformational changes, but they are random variables used by the code.
Recording $(R_i,\tau_i)$ preserves the exact random-variable density and makes
no assumption that the neural network is perfectly equivariant. When these
variables are resampled from the same priors in proposal and target, their
density factors cancel in the TPS acceptance ratio.

The relation can be summarized as follows:
\begin{itemize}[leftmargin=*]
  \item Quotient-space diffusion removes redundant movement within an
  equivalence class by construction.
  \item Boltz-2 uses SE(3) augmentation inside an ambient-coordinate sampler.
  \item TPS on the implemented Boltz-2 sampler must therefore operate on the
  extended random-variable path, not only on quotient coordinates.
  \item Observables such as ligand RMSD and ligand--pocket distance are
  computed after consistent alignment or using invariant distances, so they
  probe structural variation rather than arbitrary global motion.
\end{itemize}

\section{Supplementary Note 3: Evaluation Metrics}
\label{sec:metrics}

\subsection{Biased collective variables}

The primary OPES collective variable is
$\mathbf{s}=(s_1,s_2)$.

\paragraph{Ligand pose RMSD.}
$s_1$ is the Kabsch-aligned heavy-atom RMSD of the terminal ligand coordinates
to a reference pose:
\begin{equation}
  s_1(x)=
  \sqrt{\frac{1}{N_{\mathrm{lig}}}
  \sum_{a=1}^{N_{\mathrm{lig}}}
  \|R^\star x_a+t^\star-x^{\mathrm{ref}}_a\|^2},
\end{equation}
where $(R^\star,t^\star)$ is the optimal rigid alignment.

\paragraph{Ligand--pocket distance.}
$s_2$ is the distance between the ligand center of mass and the pocket
C$\alpha$ centroid or nearest-pocket C$\alpha$ summary used by the analysis
script:
\begin{equation}
  s_2(x)=\left\|r_{\mathrm{COM}}^{\mathrm{lig}}(x)-
  r_{\mathrm{pocket}}^{C\alpha}(x)\right\|.
\end{equation}
Together, $s_1$ distinguishes pose deviation and $s_2$ distinguishes pocket
engagement.

\subsection{Diagnostic structural metrics}

\begin{description}[leftmargin=!,labelwidth=1.55in]
  \item[Contact order]
  Fraction of long-range C$\alpha$ contacts:
  \[
    \mathrm{CO}(x)=
    \frac{|\{(i,j): |i-j|\ge 6,\ d_{ij}<8~\text{\AA}\}|}
    {|\{(i,j): |i-j|\ge 6\}|}.
  \]
  \item[Steric clashes]
  Count of non-bonded atom pairs below a distance threshold, typically
  $3.0~\text{\AA}$ for coarse screening.
  \item[lDDT]
  Alignment-free local distance difference test~\cite{mariani2013lddt}.
  High lDDT to a reference may indicate local structural agreement or
  memorization depending on the experimental design.
  \item[Ramachandran outliers]
  Fraction of residues outside allowed backbone $(\phi,\psi)$ regions.
  \item[Radius of gyration]
  $R_g=\sqrt{M^{-1}\sum_m\|x_m-\bar x\|^2}$.
  \item[Boltz pLDDT]
  Mean confidence score from the model confidence head. High pLDDT coupled
  with large geometric deviation is treated as an overconfidence warning.
\end{description}

\subsection{PoseBusters-style ligand diagnostics}

The implementation can evaluate PoseBusters-style ligand checks using either
the upstream CPU PoseBusters package or a tensor-native GPU subset. Exposed
checks include ligand RMSD thresholds, ligand--pocket distance thresholds,
contact counts, hydrogen-bond counts, clash counts, ligand maximum extent and
bounding-box volume. These are useful OPES diagnostics but were not used as
the primary two-dimensional bias in the main MEK1--FZC result.

\subsection{Ensemble-distribution metrics}

For comparing TPS-derived ensembles to molecular-dynamics or generative
ensemble baselines, the repository includes torsion Jensen--Shannon distance,
one-dimensional Wasserstein distance, per-residue RMSF correlation/RMSE,
per-atom Gaussian Wasserstein-2 distance, interaction fingerprint
Wasserstein distance, pairwise RMSD correlation, PCA Wasserstein-2,
transient contact Jaccard index and SASA correlation. These
metrics are documented in \path{docs/evaluation_metrics.md}; they are
secondary to the Brief Communication's path-sampling claim but are needed for
broader model comparison.

\section{Supplementary Note 4: Computational Workflow}
\label{sec:workflow}

\subsection{Inputs}

The main run uses PDB 7XLP, protein MEK1/MAP2K1 and ligand FZC. The pipeline
prepares the Boltz input YAML, processes structures into Boltz-compatible NPZ
topology files and initializes one accepted denoising trajectory.

\subsection{Execution pipeline}

\begin{enumerate}[leftmargin=*]
  \item Run Boltz-2 inference with random-number capture enabled. Store terminal
  coordinates and per-step random variables.
  \item Wrap the Boltz sampler as an OpenPathSampling-compatible engine. The
  snapshot includes coordinates, diffusion step index and recorded SE(3)/noise
  variables.
  \item Run TPS with forward shooting, backward shooting and global reshuffle
  moves. Save terminal frames and path metadata.
  \item Evaluate CVs on each terminal frame. The OPES run writes
  \texttt{cv\_values.json}, saved kernels and trajectory checkpoints.
  \item Update the OPES bias every \texttt{pace} steps, compute the biased
  Metropolis rule and log accepted/rejected moves.
  \item Post-process the terminal ensemble into marginal statistics,
  persistence/run-length summaries, basin assignments and structural renders.
\end{enumerate}

\subsection{Run parameters}

The one-dimensional MEK1--FZC run used ligand RMSD $s_1$ as the OPES CV,
4,180 Monte Carlo steps and biasfactor $\gamma=10$. The two-dimensional run
used $(s_1,s_2)$, 11,590 Monte Carlo steps, 3,184 accepted moves, acceptance
rate $27.5\%$, biasfactor $\gamma=10$, barrier $5\,k_BT$, 1,159 OPES
depositions and 82 surviving kernels after compression.

\subsection{Output artifacts}

Key artifacts are \texttt{cv\_values.json}, OPES kernel checkpoints,
\texttt{trajectory\_checkpoints/}, basin-assignment tables and rendered
terminal structures. The saved random variables allow the path probability to
be recomputed after the run.

\section{Supplementary Note 5: Algorithms}
\label{sec:algorithms}

\subsection*{Algorithm S1: TPS for explicit-noise diffusion}

\begin{enumerate}[leftmargin=*]
  \item Input: forward kernel $G$, random-variable priors
  $p(\xi_i)$, endpoint indicators $H_A,H_B$ and an initial reactive path.
  \item Select a move: forward shooting, backward shooting or global reshuffle.
  \item For forward shooting, choose $k$, keep the prefix, resample all
  $\xi_i$ for $i\ge k$ and integrate forward.
  \item For backward shooting, choose $k$, construct a new prefix using the
  backward proposal $q_{\mathrm{bwd}}$ and compute
  Eq.~\eqref{eq:delta-bwd}.
  \item For global reshuffle, draw a fresh $x_0$ and all $\xi_i$ from their
  priors and integrate the entire path.
  \item Reject immediately if endpoint indicators fail. Otherwise apply the
  relevant Metropolis probability.
  \item Save the accepted path, terminal frame, random variables and CVs.
\end{enumerate}

\subsection*{Algorithm S2: TPS+OPES}

\begin{enumerate}[leftmargin=*]
  \item Initialize the TPS path and an empty OPES kernel list.
  \item At Monte Carlo step $n$, propose a TPS move and evaluate
  $\mathbf{s}^{\mathrm{n}}=\mathbf{s}(x_L^{\mathrm{n}})$ and
  $\mathbf{s}^{\mathrm{o}}=\mathbf{s}(x_L^{\mathrm{o}})$.
  \item Compute the OPES-bias difference
  $\Delta V=V_{n-1}(\mathbf{s}^{\mathrm{n}})-V_{n-1}(\mathbf{s}^{\mathrm{o}})$.
  \item Accept with
  \[
    \HAB(\Xnew)\min\left[1,
    \exp(\Delta_{\mathrm{TPS}}-\Delta V/k_BT)\right].
  \]
  \item Every \texttt{pace} steps, deposit or merge a kernel centered at the
  accepted terminal CV.
  \item Continue until the target number of Monte Carlo steps is reached; use
  the saved OPES state for reweighted marginal summaries.
\end{enumerate}

\section{Supplementary Note 6: Secondary Results and Controls}
\label{sec:secondary}

A separate 228-residue two-chain co-folding control was run without MSA
information to test whether TPS detects a sharply funneled denoising
landscape. The system used $M=1792$ padded atoms, $L=32$ denoising steps and
initial noise scale $\sigma_0=2560~\text{\AA}$. A 100-step TPS run produced a
shooting acceptance rate of approximately 53\%. Terminal C$\alpha$ RMSD values
sampled every 10 Monte Carlo steps spanned $0.48$--$0.69~\text{\AA}$, with
mean $0.59\pm0.06~\text{\AA}$, while OpenMM-minimized energies ranged from
$-42{,}201$ to $-41{,}653$ kJ mol$^{-1}$.

This control supports the interpretation that TPS can detect when stochastic
freedom in the sampler collapses to a narrow structural attractor. It is not
used as the main Brief Communication result because the MEK1--FZC case more
directly addresses protein--ligand memorization and rare terminal basins.

\begin{table}[h]
\centering
\caption{Summary of the MEK1--FZC two-dimensional TPS+OPES run.}
\begin{tabular}{lcc}
\toprule
Quantity & Value & Note\\
\midrule
Monte Carlo steps & 11,590 & all logged proposals\\
Accepted moves & 3,184 & 27.5\% acceptance\\
OPES depositions & 1,159 & before compression\\
Surviving kernels & 82 & after compression\\
$s_1>1.0~\text{\AA}$ & 8.8\% & all logged steps\\
Mean high-RMSD dwell & $\sim 73$ steps & threshold $s_1>1.0~\text{\AA}$\\
Maximum high-RMSD dwell & 196 steps & threshold $s_1>1.0~\text{\AA}$\\
\bottomrule
\end{tabular}
\end{table}

\section*{References}

\begin{thebibliography}{99}

\bibitem{karras2022}
T.~Karras, M.~Aittala, T.~Aila and S.~Laine,
\textit{Adv.~Neural Inf.~Process.~Syst.} \textbf{35}, 26565 (2022).

\bibitem{xu2026}
Y.~Xu, Y.~Wang, S.~Luo, K.~Gao, T.~He, D.~He and C.~Liu,
\textit{ICLR} (2026), arXiv:2604.21809.

\bibitem{mariani2013lddt}
V.~Mariani, M.~Biasini, A.~Barbato and T.~Schwede,
\textit{Bioinformatics} \textbf{29}, 2722 (2013).

\end{thebibliography}

\end{document}
